% !TEX root = ../main.tex
\Cref{ch:machine} prints the main parts of \texttt{program/erdos915.py}.
This appendix prints the remaining parts, in the order they appear in the file, and the output of two full runs, one with each heuristic search.

\section{The remaining parts of the program}\label{sec:program-rest}

The file starts with a short description and the imports.

\programpart{header}{The description and the imports.}

Two helpers list the vertex pairs that must be checked, ordered in directed models, and count the edges of a graph.

\programpart{pairs}{The vertex pairs and the edge count.}

The heuristic searches keep the number of disjoint routes of every pair. The exception \texttt{OutOfTime} stops a search at its time limit.

\programpart{connectivities}{The connectivity of every pair, and the time limit.}

A move builds a new graph and leaves the old one unchanged.

\programpart{apply}{Making a move.}

The file checker checks a saved file, \texttt{witnesses\_tabu.json} or \texttt{witnesses\_greedy.json}, without the rest of the program (\Cref{sec:certificates}). For each saved graph it checks that every hyperedge is one the model allows, with a whole number of copies up to the largest multiplicity, that every hyperedge lists its vertices and its tails in increasing order, so that it has only one written form, that no hyperedge is listed twice, that the copies add up to the size, that every vertex pair has exactly one certificate entry, and that every certificate is correct.

\programpart{verifier}{The file checker of the saved graphs and certificates.}

The main function first prints the versions of Python, \texttt{networkx} and \texttt{nauty}, the operating system, the number of processor cores, the time budget, the heuristic search and the date. Then it runs every case in parallel, prints one table per problem, and saves every graph with its certificate.

\programpart{main}{Running every case and printing the tables.}

\section{Output}\label{sec:output}

\Cref{lst:output} is the output of the program with tabu search, unchanged.
The command \texttt{python erdos915.py > output\_tabu.txt} in the folder \texttt{program} writes it to a file, and this thesis reads that file directly, so no step lies between the program and the page.
The first four lines of the output give the software versions, the operating system, the number of processor cores, the time budget per case, the heuristic search and the date of the run. The computer was an Apple M5 with ten processor cores.
The searches stop at time limits and run in parallel, so the unfinished searches depend on the speed and the load of the computer. A run on another computer, or a second run on the same one, can print other lower bounds, even with the same random seed.
Every entry is the number of edges, arcs or hyperedges of the densest feasible graph found, and every one of these graphs passed its certificate. The same run wrote these graphs and their certificates to \texttt{program/witnesses\_tabu.json}.
An entry marked \texttt{*} comes from a finished search, so it is the exact maximum, and no run can print a larger value for it.
An unmarked entry is a lower bound only, and a longer run may raise it.
For hypergraphs the rows start at $n=3$.
Each directed hypergraph problem has three tables, one for each version of \Cref{sec:variants}.
Wherever the searches for the forward and the backward version both finished, the two tables agree, as \Cref{lem:dir-hyper-duality} predicts.
In the general model the search finished almost only at $n=3$. Its other entries are lower bounds, and some lie below the forward entries, although every forward hypergraph also belongs to the general model.

% The file is read directly. First its first six lines (the run record and
% the legend), then every table by its line range: one row per model, edge
% separation on the left and vertex separation on the right.
{\lstinputlisting[style=outputlisting, captionpos=t, lastline=6, caption={The output of the program with tabu search, for $2\le m\le 7$ and $n\le 7$.}, label={lst:output}]{program/output_tabu.txt}}
\newcommand{\outputfile}{program/output_tabu.txt}
\newcommand{\outputtable}[2]{%
  \begin{minipage}[t]{0.47\textwidth}
  \lstinputlisting[style=outputlisting, linerange={#1-#2}, basewidth=0.43em, xleftmargin=4pt]{\outputfile}
  \end{minipage}}
\newcommand{\outputpair}[4]{\noindent\outputtable{#1}{#2}\hfill\outputtable{#3}{#4}\par\medskip}
\outputpair{8}{15}{17}{24}
\outputpair{26}{33}{35}{42}
\outputpair{44}{50}{52}{58}
\outputpair{60}{66}{68}{74}
\outputpair{76}{83}{85}{92}
\outputpair{94}{101}{103}{110}
\outputpair{112}{118}{120}{126}
\outputpair{128}{134}{136}{142}
\outputpair{144}{150}{152}{158}
\outputpair{160}{166}{168}{174}
\outputpair{176}{182}{184}{190}
\outputpair{192}{198}{200}{206}

\section{Output with random greedy growth}\label{sec:output-greedy}
\Cref{lst:output-greedy} is the output of the program with random greedy growth.
The command \texttt{python erdos915.py greedy > output\_greedy.txt} writes it, on the same computer and with the same time budget as \Cref{lst:output}. The same run wrote its graphs and certificates to \texttt{program/witnesses\_greedy.json}.

\renewcommand{\outputfile}{program/output_greedy.txt}
{\lstinputlisting[style=outputlisting, captionpos=t, lastline=6, caption={The output of the program with random greedy growth, for $2\le m\le 7$ and $n\le 7$.}, label={lst:output-greedy}]{\outputfile}}
\outputpair{8}{15}{17}{24}
\outputpair{26}{33}{35}{42}
\outputpair{44}{50}{52}{58}
\outputpair{60}{66}{68}{74}
\outputpair{76}{83}{85}{92}
\outputpair{94}{101}{103}{110}
\outputpair{112}{118}{120}{126}
\outputpair{128}{134}{136}{142}
\outputpair{144}{150}{152}{158}
\outputpair{160}{166}{168}{174}
\outputpair{176}{182}{184}{190}
\outputpair{192}{198}{200}{206}
